\begin{definition}[Alexander-dual coordinate monomials]
\label{def:philosophy-alexander-dual-coordinate-monomials}
Let $R_{0}=k[x_{p}:p\in E]$ be the squarefree polynomial ring on the
external codons. For each trigger coordinate $i\in\Omega$, define
$$
  U_{i}=\prod_{\sigma(p)\ni i}x_{p}.
$$
Since a fixed coordinate lies in $2^{3}$ quotient supports and each
support has two $f_{3}$-lifts, every $U_{i}$ has degree $16$.
\end{definition}

\begin{theorem}[Alexander-dual blocker ideal]
\label{thm:philosophy-alexander-dual-blocker-ideal}
Let $\mathcal{B}$ be the blocker complex. Its Alexander-dual ideal is
$$
  I_{\mathcal{B}}^{\vee}
  =
  \langle
  U_{s_{1}},U_{f_{1}},U_{s_{2}},U_{f_{2}}
  \rangle.
$$
\end{theorem}
\begin{proof}
The facets of $\mathcal{B}$ are the four maximal blockers
$B_{\neg i}=\{p\in E:i\notin\sigma(p)\}$. Alexander duality sends each
facet to its complement in $E$:
$$
  E\setminus B_{\neg i}=\{p\in E:i\in\sigma(p)\}.
$$
The squarefree monomial on this complement is precisely $U_{i}$.
\end{proof}

\begin{theorem}[LCM degrees]
\label{thm:philosophy-alexander-dual-lcm-degrees}
For $S\subseteq\Omega$ with $|S|=k$,
$$
  \deg\operatorname{lcm}(U_{i}:i\in S)
  =
  32-2^{5-k}.
$$
Thus the lcm degrees are
$$
\begin{array}{c|r|r}
  |S| & \text{number of subsets} & \text{lcm degree}\\
  \hline
  1 & 4 & 16\\
  2 & 6 & 24\\
  3 & 4 & 28\\
  4 & 1 & 30.
\end{array}
$$
\end{theorem}
\begin{proof}
The lcm contains exactly those external codons whose support intersects
$S$. Among the $15$ nonempty quotient supports, the supports disjoint
from $S$ are the nonempty subsets of $\Omega\setminus S$, hence number
$2^{4-k}-1$. Therefore the quotient supports meeting $S$ number
$15-(2^{4-k}-1)=16-2^{4-k}$, and each has two $f_{3}$-lifts.
\end{proof}

\begin{theorem}[Minimal Taylor resolution]
\label{thm:philosophy-alexander-dual-minimal-taylor-resolution}
For
$$
  J=I_{\mathcal{B}}^{\vee},
$$
the Taylor resolution is minimal, with graded shifts
$$
  0
  \to
  R_{0}(-30)
  \to
  R_{0}(-28)^{4}
  \to
  R_{0}(-24)^{6}
  \to
  R_{0}(-16)^{4}
  \to
  J
  \to
  0.
$$
Equivalently, the quotient $R_{0}/J$ has the resolution
$$
  0
  \to
  R_{0}(-30)
  \to
  R_{0}(-28)^{4}
  \to
  R_{0}(-24)^{6}
  \to
  R_{0}(-16)^{4}
  \to
  R_{0}
  \to
  R_{0}/J
  \to
  0.
$$
\end{theorem}
\begin{proof}
The four generators have lcms whose degrees depend only on the size of
the chosen coordinate subset and strictly increase with that size.
Hence no Taylor differential can cancel equal multidegrees. The ranks
of the Taylor modules are $\binom{4}{1},\binom{4}{2},\binom{4}{3}$ and
$\binom{4}{4}$, and the shifts are those of
\autoref{thm:philosophy-alexander-dual-lcm-degrees}.
\end{proof}

\begin{corollary}[Dual compression of the trigger hypergraph]
\label{cor:philosophy-dual-compression-trigger-hypergraph}
The Stanley-Reisner ideal $I_{\mathcal{B}}$ has $294$ minimal
squarefree generators, but its Alexander dual has only four. The Betti
rank pattern of the dual resolution is
$$
  1,\ 4,\ 6,\ 4,\ 1.
$$
\end{corollary}
\begin{proof}
The $294$ generators of $I_{\mathcal{B}}$ are the minimal full triggers
from \autoref{thm:philosophy-stanley-reisner-trigger-ideal}. The
dual ideal has the four generators of
\autoref{thm:philosophy-alexander-dual-blocker-ideal}; its Taylor ranks
are the binomial coefficients in four variables.
\end{proof}

\begin{theorem}[Reliability from the lcm lattice]
\label{thm:philosophy-reliability-from-lcm-lattice}
For $q=1-p$, the non-trigger reliability polynomial is
$$
  P_{\mathrm{non}}(p)
  =
  \sum_{\emptyset\neq S\subseteq\Omega}
  (-1)^{|S|+1}
  q^{\deg\operatorname{lcm}(U_{i}:i\in S)}
  =
  4q^{16}-6q^{24}+4q^{28}-q^{30}.
$$
\end{theorem}
\begin{proof}
The event $U_{i}$ is absent is the event that coordinate $i$ is missed.
Inclusion-exclusion over the four missed-coordinate events uses exactly
the lcm degrees of \autoref{thm:philosophy-alexander-dual-lcm-degrees}.
The coefficients are $4,-6,4,-1$ according to subset size.
\end{proof}

\begin{theorem}[Alexander-dual trigger algebra theorem]
\label{thm:philosophy-alexander-dual-trigger-algebra-theorem}
Beyond $M^{\star}=\mathrm{NNR}\cup\mathrm{CUY}$, the trigger side has
$294$ codon-level minimal full triggers, while the blocker-dual side is
the four-generator squarefree monomial ideal
$$
  I_{\mathcal{B}}^{\vee}
  =
  \langle
  U_{s_{1}},U_{f_{1}},U_{s_{2}},U_{f_{2}}
  \rangle.
$$
Its lcm degrees are $16,24,28,30$, its minimal Taylor ranks are
$1,4,6,4,1$, and its lcm-lattice inclusion-exclusion is exactly
$$
  4q^{16}-6q^{24}+4q^{28}-q^{30}.
$$
\end{theorem}
\begin{proof}
The four-generator description is
\autoref{thm:philosophy-alexander-dual-blocker-ideal}. The shifts and
Betti pattern are
\autoref{thm:philosophy-alexander-dual-minimal-taylor-resolution}. The
probability polynomial is
\autoref{thm:philosophy-reliability-from-lcm-lattice}.
\end{proof}
